\usepackage[a4paper,top=20mm,bottom=22mm,left=19mm,right=19mm,headheight=25pt]{geometry}
\usepackage[table]{xcolor}
\usepackage{amsmath,amssymb}
% Native Unicode math retains searchable/copyable operators and extensible
% delimiters; legacy CMSY/CMEX fonts lacked ToUnicode tables in XeTeX output.
\usepackage{unicode-math}
\setmathfont{latinmodern-math.otf}
\renewcommand{\boldsymbol}[1]{\symbfit{#1}}
\usepackage{siunitx}
\usepackage{booktabs,tabularx,array,multicol,longtable}
\usepackage{enumitem}
\usepackage{graphicx,float,caption}
\usepackage{tikz,pgfplots}
\pgfplotsset{compat=1.18}
\usepgfplotslibrary{polar}
\usetikzlibrary{arrows.meta,positioning,calc,angles,quotes,patterns,
  decorations.pathmorphing,decorations.markings,intersections,
  shapes.geometric,backgrounds,3d,fit,spy}
\usepackage[most]{tcolorbox}
\usepackage{fancyhdr}
\usepackage{hyperref,bookmark}
\usepackage{xparse}

\definecolor{Ink}{HTML}{202326}
\definecolor{DeepBlue}{HTML}{294E69}
\definecolor{Blue}{HTML}{47687D}
\definecolor{PaleBlue}{HTML}{F3F5F6}
\definecolor{PaleGold}{HTML}{F6F6F5}
\definecolor{RuleGray}{HTML}{D3D7DA}
\definecolor{TextGray}{HTML}{61686D}
\definecolor{SoftGray}{HTML}{F6F6F6}
\color{Ink}

\sisetup{per-mode=symbol,inter-unit-product=\ensuremath{{}\cdot{}},exponent-product=\times}
\hypersetup{
  colorlinks=true,linkcolor=DeepBlue,urlcolor=Blue,bookmarksdepth=2,
  pdfauthor={物理竞赛二轮资料编写组}
}

\setlength{\parindent}{2em}
\setlength{\parskip}{0.18em}
\setlist{nosep,leftmargin=2em}
\setcounter{secnumdepth}{2}
\setcounter{tocdepth}{1}
\allowdisplaybreaks[2]
\emergencystretch=2em

\ctexset{
  chapter={format=\raggedright\LARGE\bfseries\color{Ink},beforeskip=0pt,afterskip=18pt},
  section={format=\Large\bfseries\color{DeepBlue},beforeskip=1.8ex,afterskip=1.0ex},
  subsection={format=\large\bfseries\color{Ink},beforeskip=1.5ex,afterskip=0.8ex},
  subsubsection={format=\normalsize\bfseries\color{Ink}}
}

\pagestyle{fancy}
\fancyhf{}
\fancyhead[L]{\small\color{TextGray}高中物理竞赛二轮讲义}
\fancyhead[R]{\small\color{TextGray}\nouppercase{\leftmark}}
\fancyfoot[C]{\small\color{TextGray}\thepage}
\renewcommand{\headrulewidth}{0.4pt}
\renewcommand{\headrule}{\hbox to\headwidth{\color{RuleGray}\leaders\hrule height \headrulewidth\hfill}}

\newcommand{\dd}{\mathop{}\!\mathrm d}
\newcommand{\e}{\mathrm e}
\newcommand{\ii}{\mathrm i}
\newcommand{\vect}[1]{\symbfit{#1}}
\newcommand{\uvec}[1]{\hat{\vect{#1}}}
\newcommand{\order}[1]{\mathcal O\!\left(#1\right)}
\DeclareMathOperator{\proj}{proj}

\renewcommand{\theequation}{\thechapter.\arabic{equation}}
\renewcommand{\thefigure}{\thechapter.\arabic{figure}}
\renewcommand{\thetable}{\thechapter.\arabic{table}}
\makeatletter
\@addtoreset{equation}{chapter}
\@addtoreset{figure}{chapter}
\@addtoreset{table}{chapter}
\makeatother

\newcounter{workedexample}[chapter]
\renewcommand{\theworkedexample}{\thechapter.\arabic{workedexample}}
\NewDocumentEnvironment{examplebox}{m}{%
  \refstepcounter{workedexample}%
  \begin{tcolorbox}[enhanced,breakable,colback=white,colframe=RuleGray,
    boxrule=0.5pt,arc=0mm,title={例~\theworkedexample\quad #1},fonttitle=\bfseries,
    coltitle=Ink,colbacktitle=PaleBlue,borderline west={1.4pt}{0pt}{DeepBlue}]
}{\end{tcolorbox}}

\newcounter{practice}[chapter]
\renewcommand{\thepractice}{\thechapter.\arabic{practice}}
\NewDocumentEnvironment{practicebox}{m}{%
  \refstepcounter{practice}%
  \begin{tcolorbox}[enhanced,breakable,colback=white,colframe=RuleGray,
    boxrule=0.5pt,arc=0mm,title={练~\thepractice\quad #1},fonttitle=\bfseries,
    coltitle=Ink,colbacktitle=PaleGold,borderline west={1.4pt}{0pt}{TextGray}]
}{\end{tcolorbox}}

% 高阶深化卷采用 H1--H6 独立编号，不占用正文第1--28章的章号。
\newcommand{\advancedunit}[2]{%
  \clearpage
  \chapter*{#1\quad #2}%
  \addcontentsline{toc}{chapter}{#1\quad #2}%
  \markboth{#1\quad #2}{#1\quad #2}%
  \setcounter{section}{0}\setcounter{subsection}{0}%
  \setcounter{equation}{0}\setcounter{figure}{0}\setcounter{table}{0}%
  \setcounter{workedexample}{0}\setcounter{practice}{0}%
  \renewcommand{\theequation}{#1.\arabic{equation}}%
  \renewcommand{\thefigure}{#1.\arabic{figure}}%
  \renewcommand{\thetable}{#1.\arabic{table}}%
  \renewcommand{\theworkedexample}{#1.\arabic{workedexample}}%
  \renewcommand{\thepractice}{#1.\arabic{practice}}%
}

\NewDocumentEnvironment{solution}{+b}{%
  \par\smallskip\begingroup\small\noindent
  {\bfseries\color{DeepBlue}解答}\quad #1\par\endgroup\smallskip
}{}

\newtcolorbox{definitionbox}[1]{enhanced,breakable,colback=PaleGold,
  colframe=RuleGray,boxrule=0.5pt,arc=0mm,title={#1},fonttitle=\bfseries,
  coltitle=Ink,colbacktitle=PaleGold,borderline west={1.2pt}{0pt}{TextGray}}

\newcommand{\bookcaption}[1]{%
  \refstepcounter{figure}%
  \par\nopagebreak[4]\smallskip
  {\centering\small\color{TextGray}图~\thefigure：#1\par}\smallskip}

\tikzset{
  every node/.append style={fill=none,inner sep=1.1pt},
  labelclean/.style={fill=none,inner sep=0.7pt,outer sep=1.2pt,text=Ink},
  point/.style={circle,fill=Ink,draw=none,inner sep=1.6pt},
  vector/.style={-{Latex[length=2.5mm]},very thick,draw=DeepBlue},
  component/.style={-{Latex[length=2.2mm]},thick,draw=TextGray},
  axis/.style={-{Latex[length=2mm]},thin,draw=TextGray},
  guide/.style={densely dashed,draw=TextGray},
  curve/.style={very thick,draw=DeepBlue},
  ground/.style={very thick,draw=Ink},
  spring/.style={decorate,decoration={zigzag,segment length=4pt,amplitude=2pt},thick,draw=DeepBlue}
}
